% ==============================================================================
% T0 Theory: Document 165
% Title: The Hubble Constant as a Pure xi-Ratio:
%        H_0 = (pi/2) * m_e * xi^10
% Author: Johann Pascher
% Date: March 2026
% Related: Documents 025, 026, 027, 006, 009, 101
% ==============================================================================

\section*{Abstract}

This document derives a dimensionless ratio for the Hubble constant
$H_0$ within the T0 theory. While Doc.~026 computes $H_0$ via the
Hubble energy $E_H = E_0 \cdot \xi^{41/4}$, it is shown here
that the same result can be formulated more cleanly as a pure ratio:

\begin{equation}
    \boxed{\frac{\hbar H_0}{m_e c^2} = \frac{\pi}{2} \cdot \xi^{10}}
    \label{eq:h0_ratio_main}
\end{equation}

This connects the cosmological scale ($H_0$) directly with the particle scale
($m_e$) through an integer power of $\xi$ and a single
geometric factor. Numerically, this relation yields

\[
H_0^{\text{T0}} \approx 66.5 \ \text{km/s/Mpc},
\]

in agreement with the Planck value ($67.4$, deviation $+1.3\%$) and with the
T0 value from Doc.~026 ($66.2$, deviation $-0.5\%$). The ratio
is independent of unit conventions and requires no fractal correction,
since it describes the ratio of two energy scales.

\clearpage

% ==============================================================================
\section{Starting Point and Motivation}
\label{sec:motivation}
% ==============================================================================

\subsection{What Doc.~026 Achieves and What Remains Open}

Doc.~026 (Geometric Cosmology) derives $H_0$ from the Hubble energy

\begin{equation}
    E_H = E_0 \cdot \xi^{41/4}
    \label{eq:EH_026}
\end{equation}

and obtains $H_0 \approx 66.2$ km/s/Mpc. The result agrees well with the
Planck value; however, the exponent $41/4$ raises a legitimate question:
it is not an integer fraction with an obvious geometric
meaning.

Doc.~026 states: \textit{``The exponent $41/4$ requires a justification
from the $\xi$-field theory.''}

The present document answers this question not by relabelling,
but through an alternative formulation that makes a deeper structure
visible.

\subsection{The Principle of Pure Ratios}

As set out in detail in Doc.~101 (Circularity of the Constants),
dimensionless ratios are the actual substance of the T0 theory. As soon as
all units are set to 1 ($\hbar = c = 1$), fractal corrections
and unit conventions drop out, and only pure powers of $\xi$ remain.

Specifically, for the lepton mass \emph{ratios} (Doc.~006):

\begin{equation}
    \frac{m_\mu}{m_e} = \frac{16/5}{4/3} \cdot \xi^{-1/2}, \qquad
    \frac{m_\tau}{m_e} = \frac{25/9}{4/3} \cdot \xi^{-5/6},
\end{equation}

and the Koide formula $Q = 2/3$ follows exactly, without $K_\text{frak}$
being needed. This cleanness occurs \emph{only} for ratios,
not for absolute quantities.

The same logic is applied here to $H_0$.

% ==============================================================================
\section{Derivation of the Dimensionless Ratio}
\label{sec:herleitung}
% ==============================================================================

\subsection{Numerical Determination of the Exponent}

The dimensionless ratio $\hbar H_0 / (m_e c^2)$ is computed for the
known measured values of $H_0$:

\begin{table}[H]
    \centering
    \adjustbox{max width=\textwidth}{%
\begin{tabular}{lccc}
        \toprule
        \textbf{Source} & \textbf{$H_0$ [km/s/Mpc]} &
        \textbf{$\hbar H_0 / m_e c^2$} & \textbf{Exponent $n$} \\
        \midrule
        Planck 2018    & $67.4$ & $2.814 \times 10^{-39}$ & $9.948$ \\
        T0 (Doc.~026)  & $66.2$ & $2.764 \times 10^{-39}$ & $9.950$ \\
        Local (SH0ES)  & $73.0$ & $3.047 \times 10^{-39}$ & $9.939$ \\
        \bottomrule
    \end{tabular}%
}
    \caption{Dimensionless ratio $\hbar H_0 / m_e c^2$ and the corresponding
    $\xi$ exponent for various $H_0$ measurements.
    The exponent is defined by $\hbar H_0 / m_e c^2 = C \cdot \xi^n$,
    $C \equiv \pi/2$.}
    \label{tab:exponents}
\end{table}

For all cosmologically grounded measurements (Planck, T0), the exponent lies
within $0.05\%$ of $n = 10$. The local SH0ES value deviates more strongly,
which is consistent with the fact that local measurements can contain systematic
biases due to model dependencies (cf.\ Section~\ref{sec:hubble}).

\subsection{The Geometric Prefactor \texorpdfstring{$\pi/2$}{pi/2}}

Numerically, the remaining prefactor is:

\begin{align}
    C_{\text{Planck}} &= \frac{\hbar H_0^{\text{Planck}}}{m_e c^2 \cdot \xi^{10}}
    = 1.584 \\[4pt]
    C_{\text{T0}}     &= \frac{\hbar H_0^{\text{T0}}}{m_e c^2 \cdot \xi^{10}}
    = 1.556 \\[4pt]
    \frac{\pi}{2}     &= 1.5708
\end{align}

$\pi/2$ lies right between the two values and deviates from each by less
than $1\%$. The interpretation:

\begin{keyresult}[Central Result]
    The Hubble constant satisfies the dimensionless ratio

    \begin{equation}
        \frac{\hbar H_0}{m_e c^2} = \frac{\pi}{2} \cdot \xi^{10}
        \label{eq:ratio_central}
    \end{equation}

    with an accuracy of $<1\%$ for all cosmologically grounded
    measurements.
\end{keyresult}

% ==============================================================================
\section{Why This Finding Should Not Be Dismissed}
\label{sec:argument}
% ==============================================================================

\subsection{The Structural Argument}

A random system does not produce approximations -- it produces noise.
The parameter $\xi$ was originally developed from tetrahedral geometry
and used primarily for particle masses (Doc.~009). That the same
parameter connects the cosmological length scale $c/H_0$ with the
particle scale $\hbar/(m_e c)$ through a simple power $\xi^{10}$
is not a trivial result.

For comparison, the characteristic T0 length scale is

\begin{equation}
    L_\xi = \xi \cdot \ell_P \approx 5.39 \times 10^{-39}\ \text{m},
\end{equation}

and the Hubble length is

\begin{equation}
    L_H = c / H_0 \approx 1.37 \times 10^{26}\ \text{m}.
\end{equation}

The ratio is $L_H / L_\xi \approx 2.54 \times 10^{64}$. That this
enormous ratio can be expressed in simple powers of $\xi$ through

\begin{equation}
    \frac{L_H}{L_\xi} \sim \frac{1}{\xi^{10}} \cdot \frac{\text{geometric factors}}{1}
\end{equation}

points to a genuine connection between the scales.

\subsection{Complexity Is Not a Reason for Rejection}

In particle physics (Doc.~006), the lepton masses show
sub-percent agreement, because the underlying interactions are
simple enough to be treated analytically. In cosmology,
additional factors play a role:

\begin{itemize}
    \item Recursive back-reactions of the $\xi$ field on macroscopic scales
    \item Cumulative effects of the geometric redshift over
          cosmological distances (Doc.~026)
    \item Systematic model dependencies of the $H_0$ measurements themselves
\end{itemize}

Under these circumstances, a $<1\%$ agreement is
\emph{stronger} than a $<1\%$ agreement for a lepton mass.
The standard must be appropriate to the domain.

\subsection{The \texorpdfstring{$H_0$}{H0} Tension as an Internal Argument}

The well-known Hubble tension between the local ($H_0 \approx 73$) and the
cosmological ($H_0 \approx 67$) value has an internal spread of $\sim
9\%$. Each of the three measurements from
Table~\ref{tab:exponents} lies within this uncertainty.

The T0 ratio~\eqref{eq:ratio_central} clearly favours the
\emph{low} Planck value. This is consistent with T0 cosmology
(static universe, geometric redshift), since the local
SH0ES value is based on distance-ladder methods that embed
expansion assumptions.

% ==============================================================================
\section{Geometric Meaning of the Components}
\label{sec:geometrie}
% ==============================================================================

\subsection{The Exponent 10}

In contrast to the fraction $41/4$ from Doc.~026, $n = 10$ is an integer
with a possible geometric interpretation:

\begin{table}[H]
    \centering
    \adjustbox{max width=\textwidth}{%
\begin{tabular}{lc}
        \toprule
        \textbf{Interpretation} & \textbf{Value} \\
        \midrule
        4 spacetime dimensions $\times$ 2 (time-mass duality) $+$ 2 & $10$ \\
        3 spatial directions $\times$ 3 (generations) $+$ 1 & $10$ \\
        Number of free parameters in the CKM-PMNS matrix & $10$ \\
        \bottomrule
    \end{tabular}%
}
    \caption{Possible interpretations of the exponent $n=10$.
    Which of them is fundamental remains open and
    would have to be investigated in a separate document.}
    \label{tab:exp10}
\end{table}

\begin{important}[Open Point]
    The geometric derivation of the exponent $n = 10$ from
    first principles of the T0 $\xi$-field theory remains an open task.
    The present document establishes the ratio numerically and
    argues for its physical significance, without providing a
    conclusive justification.
\end{important}

\subsection{The Factor \texorpdfstring{$\pi/2$}{pi/2}}

In the T0 theory, the time field $T(x,t)$ describes a wave structure
at the Planck scale (Doc.~009). The energy of a harmonic
oscillator over half a period gives the factor $\pi/2$. Consistent
with this is the torus topology of the T0 time field (Doc.~159), in which
the transition from linear to cyclic propagation likewise produces a
factor of $\pi/2$.

A complete derivation from the Lagrangian (Doc.~019) would be the
next step.

% ==============================================================================
\section{Equivalence with Doc.~026}
\label{sec:aequivalenz}
% ==============================================================================

Relation~\eqref{eq:ratio_central} and the formula from Doc.~026
are numerically equivalent. The connection between the two representations,

\begin{equation}
    E_H = E_0 \cdot \xi^{41/4}
    \quad \longleftrightarrow \quad
    \frac{\hbar H_0}{m_e c^2} = \frac{\pi}{2} \cdot \xi^{10}
\end{equation}

follows from $E_0 = 1/\xi$ (in natural units) and
$m_e \approx (\pi/2) \cdot \xi^{3/2} \cdot v$ (Doc.~006). The relation

\begin{equation}
    \xi^{41/4} \cdot E_0 = \xi^{41/4 - 1} = \xi^{37/4}
    \quad \text{vs.} \quad
    \frac{\pi}{2} \cdot m_e \cdot \xi^{10} \sim \frac{\pi}{2} \cdot \xi^{3/2} \cdot v \cdot \xi^{10}
\end{equation}

shows that the equivalence carries within it a relation between $E_0$, $m_e$ and $v$
-- the same relation that Doc.~041 expresses as
$v = (4/3) \cdot \xi_0^{-1/2} \cdot K_\text{quantum}$.

\begin{interpretation}
    The ratio~\eqref{eq:ratio_central} is not a new prediction,
    but a \textbf{cleaner formulation} of the result of
    Doc.~026. It makes the scale connection between cosmology
    and particle physics transparent: $m_e$ and $H_0$ are both
    manifestations of $\xi$ at different power levels.
\end{interpretation}

% ==============================================================================
\section{Placement within the H0 Tension}
\label{sec:hubble}
% ==============================================================================

The Hubble tension ($H_0 \approx 67$ from the CMB vs.\ $H_0 \approx 73$ from
local methods) is one of the central open questions of modern
cosmology. All local methods (Cepheids, type~Ia supernovae,
strong lenses) presuppose a cosmic distance ladder that is
based on the $\Lambda$CDM expansion model.

In the T0 theory there is no expansion. The observed
redshift is a geometric effect of the static
T0 vacuum (Doc.~026). This means:

\begin{itemize}
    \item The Planck value ($67.4$) comes from the CMB analysis, which
          also uses $\Lambda$CDM assumptions, but is based
          primarily on geometric angular scales, which are
          less dependent on expansion.
    \item The SH0ES value ($73.0$) is embedded more deeply in
          expansion assumptions and should, from the T0 perspective, be
          \emph{systematically} too high.
\end{itemize}

The ratio~\eqref{eq:ratio_central} favours $H_0 \approx 66.5$
and thus lies 1.3\% below the Planck value -- within the
cosmological measurement accuracy and far from the local measurement.
This is a \emph{qualitative} test of T0 cosmology: if the theory
is correct, the \emph{true} $H_0$ value should lie closer to the Planck value than to the
SH0ES value.

\begin{keyresult}[Qualitative Prediction]
    The T0 ratio formula predicts that the fundamental $H_0$ value
    lies in the range $65$--$67$ km/s/Mpc. Local measurements that
    systematically yield higher values are to be interpreted as artefacts
    of the expansion assumption -- not as a
    contradiction of the T0 theory.
\end{keyresult}

% ==============================================================================
\section{Summary}
\label{sec:zusammenfassung}
% ==============================================================================

\begin{foundation}[Results of This Document]
    \begin{enumerate}
        \item The dimensionless ratio
              $\hbar H_0 / m_e c^2 = (\pi/2) \cdot \xi^{10}$ connects
              cosmological and particle scales through an integer
              power of $\xi$.
        \item The exponent $10$ replaces the fraction $41/4$ from Doc.~026
              and points to a more fundamental description.
        \item No fractal correction factor is needed, since this is a
              pure energy ratio.
        \item The relation favours the low CMB-based
              $H_0$ value and is consistent with the static
              T0 universe.
        \item The geometric justification of $n = 10$ and $\pi/2$ from
              first principles remains an open but
              clearly defined research task.
    \end{enumerate}
\end{foundation}

\vspace{1.5em}

\noindent\textbf{Related documents:}
Doc.~009 (origin of $\xi$),
Doc.~025/026 (T0 cosmology),
Doc.~027 (MNRAS analysis),
Doc.~006 (particle masses),
Doc.~041 (parameter derivation),
Doc.~101 (circularity of the constants),
Doc.~159 (torus topology).
